\documentclass[11pt]{article}

% Course rules: ACL template, 2-column, max 4 pages body.
% References and AI disclosure do NOT count toward the limit.
% "final" = non-anonymous camera-ready look (no line numbers).
\usepackage[final]{acl}

\usepackage{times}
\usepackage{latexsym}
\usepackage[T1]{fontenc}
\usepackage[utf8]{inputenc}
\usepackage{microtype}
\usepackage{inconsolata}
\usepackage{graphicx}
\usepackage{booktabs}

% Placeholder title in Title Case. Eridho can replace it.
% The dropped question ("when does tool use pay off?") stays out of the title.
\title{How Tool Use Affects Accuracy in ReAct and Toolformer}

\author{Eridho Rollian \\
  Universität Hamburg \\
  Seminar Understanding Large Language Models (WiSe 2026/27), Topic C12}

\begin{document}
\maketitle

\begin{abstract}
% Living draft, printed before the introduction. Revise it as the question,
% comparison, and finding settle; do a last pass when the discussion is stable.
% ACL: one paragraph, at most 200 words, thesis and conclusions.
% Cover: problem, what you compare, main finding.
TODO
\end{abstract}

\section{Introduction}
\label{sec:intro}
% Background/problem -> objective -> research question -> contribution/outline.
% Research question: how tool use affects accuracy in ReAct and Toolformer, and
% how those benefits relate to call frequency and the number of samples.
% Tokens, latency, and money are unmeasured. Do not ask whether the cost is "justified".
TODO

\section{Background}
\label{sec:background}
% Keep short (~0.4 page). In-context prompting, chain-of-thought, what "tool use" means.
TODO

\section{ReAct: Prompted Tool Use}
\label{sec:react}
% Reasoning traces interleaved with actions \citep{yao2023react}.
% Table 1 (PaLM-540B): ReAct beats chain-of-thought on FEVER (60.9 vs. 56.3
% accuracy) and trails it on HotpotQA (27.4 vs. 29.4 exact match).
% Best combination order depends on the task: HotpotQA ReAct then CoT-SC
% 35.1 exact match; FEVER CoT-SC then ReAct 64.6 accuracy.
% Figure 2 is sample efficiency (3--5 samples match CoT-SC's 21), not total cost.
% It counts CoT-SC samples only, not the combinations' ReAct steps or Wikipedia calls.
% Table 2 (HotpotQA, human labels on 50 successes + 50 failures per method):
% 6% of successful ReAct trajectories contain a hallucinated thought or fact.
TODO

\section{Toolformer: Learned Tool Use}
\label{sec:toolformer}
% Self-supervised API-call annotation \citep{schick2023toolformer}.
% Loss filter: keep a call only if it reduces LM loss. Not a correctness check
% and not a cost measure. Evaluation allows at most one call per input.
TODO

\section{Discussion}
\label{sec:discussion}
% Contribution: task-dependent gains, Table 9 reanalysis, selective calling,
% and the missing fix/break table. Do not claim cost-effectiveness.
Table~9 of \citet{schick2023toolformer} shows diminishing returns on T-REx:
raising $k$ from 3 to 10 increases the share of examples with a call by
15.3 percentage points (82.8\,\% to 98.1\,\%) but overall accuracy by only
0.6 points (52.9 to 53.5). At the same time, accuracy falls both on examples
with a call (58.0 to 54.0) and on those without one (29.0 to 22.5). This is a
composition effect, not evidence that calls became less effective: overall
accuracy is the weighted average of the two groups, and as $k$ rises,
questions move from the uncalled to the called group, which changes both
averages. By my own estimate, the newly called 15.3\,\% of questions score
about 30\,\% without a call and about 32\,\% with one, a gain of roughly
2.5 points (about 1.7 to 3.4 given the rounding of the reported figures);
since they score well below the questions already called at $k=3$, they
appear to be harder ones. This estimate rests on two assumptions that
Table~9 does not confirm: that the questions called at $k=3$ are still
called at $k=10$ and still score 58.0, and that the 1.9\,\% never called
already scored 22.5 at $k=3$. It is also a net figure, so it cannot
separate questions that a call fixes from those it breaks.

Both papers make tool use selective, but they place the decision
differently. In ReAct, two fixed rules decide when to switch between
internal reasoning and external search: ReAct falls back to CoT-SC if it
returns no answer within 7 steps on HotpotQA or 5 on FEVER, and CoT-SC
falls back to ReAct if its majority answer occurs fewer than $n/2$ times
\citep{yao2023react}. In Toolformer, the model itself decides whether to
call, through the probability it assigns to the call token, and at $k=1$
this decision appears calibrated: on T-REx it scores 44.3 on the
59.7\,\% of examples it does not call, compared with 34.9 on all examples
when calls are disabled, so it tends to call where it would do badly
without a call \citep{schick2023toolformer}. This calibration is lost at
$k=3$, where accuracy on the 17.2\,\% left uncalled falls to 29.0, below
that baseline; because these figures cover different sets of examples,
they show a pattern, not the effect of a call on any single question.
% optional:
Here, ``selective'' refers to accuracy; neither paper measures what a
call costs.

\section{Conclusion}
\label{sec:conclusion}
% Main messages, implications, limitations.
TODO

% ---- End of 4-page body. Everything below does not count. ----
\label{endofbody}

\section*{AI Disclosure}
% Required. Generated from notes/ai-log.md: which tool, which task, which section/figure.
TODO

\nocite{*} % lists all refs.bib entries; remove once you cite with \citet/\citep
\bibliography{refs}

\end{document}
